\chapter{System Implementation}
\label{chap:system}

This chapter describes how the prototype is built. It opens with
the requirements and the overall architecture, walks through each
subsystem in the order data flows, and closes with the evaluation
harness and the measured results across a nine-configuration
sweep. Code excerpts are short and chosen to show a specific
implementation choice; the full source is in the repository.
Cross-references back to Chapter~\ref{chap:foundations} appear
where an implementation choice resolves an abstract design from
the theory chapter.

\section{Requirements and Architecture}
\label{sec:sys-requirements}

\paragraph{Functional requirements.} The prototype has to satisfy
ten functional requirements. \textbf{F1--F2:} catalog browse
(home, category, product detail, search) and transactional flow
(cart, favorites, order history, simulated checkout).
\textbf{F3:} classical recommendation slots (``frequently bought
together'', ``similar products'', ``trending'', ``recommended
for you''). \textbf{F4--F5:} per-product behavior tracking
(dwell, scroll, viewport, review engagement) feeding a composite
readiness score and contextual chips with per-product cooldowns.
\textbf{F6--F8:} reactive-first advisor (always
user-initiated), review summarization with conflict surfacing,
and AI-assisted comparison with structured tradeoffs.
\textbf{F9--F10:} output validation (identifier-existence, claim
stripping, theme-fidelity) and evaluation instrumentation
(latency, token usage, versioned eval dataset).

\paragraph{Non-functional requirements.} The classical layer
targets p95 latency below 50~ms; the advisor end-to-end target
is p95 below 3~s, with the time-to-first-content target of 1~s
aspirational since tool-calling depth dominates inference speed.
Per-session token cost has to be measurable and within a
consumer-application budget; review summarization amortizes its
cost over the lifetime of a product. The LLM must never generate
factual product data, and behavior signals are aggregated
client-side with only summary statistics persisted.

\paragraph{Architecture overview.} The system is a three-layer
client-server application with two external dependencies: Clerk
for authentication and OpenRouter for LLM inference.
Figure~\ref{fig:architecture-overview} shows the layout.

\begin{figure}[ht]
  \centering
  \includegraphics[width=0.95\linewidth]{architecture-overview.pdf}
  \caption{High-level system architecture: Next.js client, Convex
    backend (queries, mutations, actions, crons, database), and
  external services (Clerk and OpenRouter).}
  \label{fig:architecture-overview}
\end{figure}

The client is a Next.js~16 application using React~19 and the
React Compiler. Pages render server-side where possible and
switch to client components only at interaction boundaries.
Tracking hooks, the advisor surface, and behavior aggregation run
on the client. The store catalog is read through Convex reactive
queries, so adding a favorite or seeding a new product propagates
to subscribed components without a manual refetch.

The backend is a single Convex deployment. Convex provides
queries (read-only, reactive, transactional), mutations (write,
transactional), actions (non-transactional, allowed to call
external services), crons, and a TypeScript schema. All
persistent state lives in Convex tables. The seed pipeline runs
as a Convex action, the recommendation cron runs as a Convex
action, and the advisor agent runs as a Convex action with
tool-calling enabled. Convex was chosen over a conventional
REST-plus-SQL-plus-queue stack because the prototype needs
reactive UI updates, scheduled batch flows, and a shared schema
between server and client, and Convex delivers these natively.
The tradeoff is vendor lock-in, accepted at prototype scale.

\section{Codebase and Deployment}
\label{sec:sys-codebase}

The repository is a Turborepo monorepo using bun workspaces:
\texttt{apps/web} (the Next.js frontend),
\texttt{packages/backend} (the Convex backend with
\texttt{convex/ai/} grouping the advisor, prompts, models,
review summaries, and eval integration),
\texttt{packages/ui} (shared design system),
\texttt{packages/eval} (the LLM eval harness),
\texttt{packages/env} (type-safe environment validation), and
\texttt{packages/config} (shared TypeScript configuration). The
eval harness depends on the backend, so scoring runs invoke the
same advisor action as the production client.

The application is deployed on a single VPS. Dokploy orchestrates
four containerized services behind a Traefik reverse proxy (TLS
via Let's Encrypt, hostname-based routing): the Next.js app, the
Convex API origin, the Convex HTTP-actions site, and the Convex
admin dashboard, with the Convex deployment self-hosted against a
co-located Postgres instance. Push-to-main triggers a GitHub
Actions workflow that runs \texttt{convex deploy} in parallel
with Dokploy's auto-deploy of the Next.js Docker image. The
posture is deliberately not hyperscale: the prototype runs on one
VPS and has to run convincingly end to end, not at scale.

\section{Data Layer}
\label{sec:sys-data}

The Convex schema defines twelve tables grouped logically into
four domains: catalog, transactional, behavior-and-AI, and
recommendation precomputation.
Figure~\ref{fig:er-diagram} shows the grouped
entity-relationship view.

\begin{figure}[ht]
  \centering
  \begin{tikzpicture}[
      node distance=0.6cm and 0.7cm,
      tab/.style={draw, rounded corners=2pt, align=left,
      font=\footnotesize, fill=white, minimum width=3.4cm, inner sep=4pt},
      grp/.style={draw, dashed, rounded corners=4pt, inner sep=10pt,
      font=\small\bfseries},
      rel/.style={-{Stealth[length=1.6mm]}, gray!70, thin}
    ]
    \node[tab] (cat) {\textbf{categories} \\ name, slug, parentId};
    \node[tab, right=of cat] (prod) {\textbf{products} \\ title,
    price, brand, \\ rating, tags, embedding};
    \node[tab, below=of prod] (rev) {\textbf{reviews} \\ productId,
    rating, text, \\ embedding};
    \node[tab, below=of cat] (rsum) {\textbf{reviewSummaries}
    \\ productId, positives, \\ negatives, conflicts};

    \node[grp, fit=(cat)(prod)(rev)(rsum), label={above:Catalog}] (catalog) {};

    \node[tab, right=2.2cm of prod] (cart) {\textbf{cartItems}
    \\ clerkUserId, productId, qty};
    \node[tab, below=of cart] (fav) {\textbf{favorites}
    \\ clerkUserId, productId};
    \node[tab, below=of fav] (ord) {\textbf{orders} \\ clerkUserId,
    items[], \\ total, status};

    \node[grp, fit=(cart)(fav)(ord), label={above:Transactional}] (trans) {};

    \node[tab, below=2.0cm of rsum] (beh) {\textbf{behaviorSessions}
    \\ sessionId, clerkUserId?, \\ productsViewed[]};
    \node[tab, right=of beh] (sug) {\textbf{suggestions}
    \\ sessionId, triggerProductId, \\ suggestedProducts[], status};
    \node[tab, right=of sug] (pref) {\textbf{userPreferences}
    \\ clerkUserId, favoriteCategories, \\ priceRange, favoriteBrands};

    \node[grp, fit=(beh)(sug)(pref), label={above:Behavior \& AI}] (bai) {};

    \node[tab, below=of bai.south, yshift=-0.2cm, xshift=-2.0cm]
    (cooc) {\textbf{productCoOccurrences} \\ productId,
    relatedProductId, \\ score (Jaccard)};
    \node[tab, right=of cooc] (caggr)
    {\textbf{coOccurrenceAggregation} \\ productIdA, productIdB,
    \\ count, batchId (temp)};

    \node[grp, fit=(cooc)(caggr), label={above:Recommendation
    precomputation}] (recg) {};

    \draw[rel] (cat) -- (prod);
    \draw[rel] (prod) -- (rev);
    \draw[rel] (prod) -- (rsum);
    \draw[rel] (prod) -- (cart);
    \draw[rel] (prod) -- (fav);
    \draw[rel] (prod) -- (ord);
    \draw[rel] (prod.south) to[bend right=10] (cooc.north);
    \draw[rel] (prod.south) to[bend left=5] (beh.east);
  \end{tikzpicture}
  \caption{Data model grouped by domain. Foreign keys flow from
    \texttt{products} into the transactional, behavior, and
    recommendation tables. Anonymous behavior sessions carry only
    \texttt{sessionId}; \texttt{clerkUserId} is attached for
  authenticated traffic.}
  \label{fig:er-diagram}
\end{figure}

The \texttt{products} table is denormalized: it stores
\texttt{category} as a string alongside \texttt{categoryId} and
keeps \texttt{rating} and \texttt{reviewCount} on the row. Every
product card render is a single indexed read with no joins, which
is what lets grid pages stream from Convex with no intermediate
compute. The table declares nine indices, each tied to a specific
recommendation-slot query pattern. The \texttt{reviews} table
carries a 3072-dimensional vector index over its embedding field,
which powers the theme-verification path in
Section~\ref{sec:sys-validation}.

\paragraph{Seed pipeline.} The seed runs as a Convex action with
a deterministic faker seed. About 200 products from the
DummyJSON feed are supplemented with faker.js attributes and
enriched with \texttt{useCases} and \texttt{goodFor} fields by a
batch Flash-Lite call modeled on COSMO~\cite{COSMO}. Roughly 50
user personas and 500 orders cluster each persona's history
around its preferences with a power-law spend distribution.
Reviews are AI-generated by a Flash-Lite batch call (20--30
reviews per product, four sentiment profiles, intentional
thematic overlap; about 4\,836 reviews total) so review
summarization has something to summarize.

\section{Recommendation Engine}
\label{sec:sys-recommendations}

The classical recommendation engine implements the four families
formalized in Section~\ref{sec:found-similarity} as a Stage~1
candidate generator. The pipeline lives in two files:
\texttt{convex/recommendations.ts} (Convex queries, mutations,
actions) and \texttt{convex/recommendationHelpers.ts} (pure
TypeScript algorithms with no Convex imports). Splitting the
algorithms from the runtime makes them independently testable.

\paragraph{Co-occurrence with Jaccard.} The
``frequently bought together'' slot reads precomputed Jaccard
scores from the \texttt{productCoOccurrences} table. Computation
runs as a daily cron action that reads all orders, builds the
pair counts and per-product order counts, and emits the
top-$N$ Jaccard-scored neighbors per product. At the seeded scale
(about 500 orders averaging three items each) a single Convex
action recomputes the full top-20-per-product set well under the
action timeout. A minimum-support filter of $\ge 2$
co-occurrences suppresses spurious pairs.

\paragraph{Content-based similarity.}
Listing~\ref{lst:content-similarity} shows the content-similarity
function in full. It is a weighted sum over six attribute
comparisons with the weights from
Section~\ref{sec:found-similarity}. The fact that this is a
12-line function with zero external dependencies is the point: at
the prototype's catalog scale, classical content-based similarity
is trivially cheap and is enough to fill the
``similar products'' slot on every product detail page.

\begin{listing}[ht]
\begin{minted}{typescript}
export function contentSimilarity(
  a: ProductAttrs, b: ProductAttrs, maxPriceRange: number
): number {
  return (
    0.25 * (a.category === b.category ? 1 : 0) +
    0.20 * jaccard(a.tags, b.tags) +
    0.15 * jaccard(a.useCases ?? [], b.useCases ?? []) +
    0.10 * jaccard(a.goodFor ?? [], b.goodFor ?? []) +
    0.10 * (a.brand === b.brand ? 1 : 0) +
    0.10 * priceProximity(a.price, b.price, maxPriceRange) +
    0.10 * ratingProximity(a.rating, b.rating)
  );
}
\end{minted}
  \caption{The content-similarity function. Weights sum to 1.0. The
    \texttt{useCases} and \texttt{goodFor} fields are produced
    once per product during seed via a batch Flash-Lite
    call~\cite{COSMO}, contributing structured similarity that
  pure attribute matching would miss.}
  \label{lst:content-similarity}
\end{listing}

\paragraph{Trending and personalized.} Trending uses the
exponential-decay formula with $\lambda \approx 0.0495$
(14-day half-life). The score is precomputed nightly into the
\texttt{trendingScore} field, so the slot is one indexed read at
query time. Personalized recommendations follow a switching
hybrid keyed on user purchase count, as formalized in
Section~\ref{sec:found-classical}: cold users fall through to
trending, lukewarm users get content similarity to favorites
plus category trending, warm users get a blend of co-occurrence
and content similarity. Signals are combined using Reciprocal
Rank Fusion with $k = 60$, followed by Maximal Marginal
Relevance at $\lambda = 0.7$ to force diversity. The full
hybrid is about 80 lines of straight-line code; no library is
involved.

Figure~\ref{fig:two-stage-pipeline} shows the two-stage pipeline,
connecting the candidate generator described above to the LLM
rerank step in Section~\ref{sec:sys-llm}.

\begin{figure}[ht]
  \centering
  \begin{tikzpicture}[
      box/.style={draw, rounded corners=2pt, align=center, minimum
      width=2.4cm, minimum height=0.95cm, font=\small},
      stage/.style={draw, thick, rounded corners=4pt, align=center,
        minimum width=4.2cm, minimum height=2.4cm,
      font=\small\bfseries, fill=gray!8},
      arr/.style={-{Stealth[length=2mm]}, thick},
      note/.style={font=\footnotesize\itshape, align=center}
    ]
    \node[box] (query) {Slot request \\ or user query};
    \node[stage, right=1.0cm of query] (s1) {Stage~1 \\[3pt]
      \normalfont\small Classical \\ \normalfont\small candidate
    \\ \normalfont\small generation};
    \node[stage, right=1.0cm of s1] (s2) {Stage~2 \\[3pt]
    \normalfont\small LLM rerank \\ \normalfont\small and explain};
    \node[box, right=1.0cm of s2] (out) {productIds \\ + rationale};
    \node[box, below=0.9cm of out] (hyd) {Client \\ hydration};
    \node[box, left=1.0cm of hyd] (db) {Convex \\ catalog};

    \draw[arr] (query) -- (s1);
    \draw[arr] (s1) -- node[above, font=\footnotesize] {top-$N$} (s2);
    \draw[arr] (s2) -- (out);
    \draw[arr] (out) -- (hyd);
    \draw[arr] (db) -- (hyd);
    \draw[arr, dashed] (db.north) |- ([yshift=-0.15cm]s1.south);

    \node[note, below=0.55cm of s1] {Convex query \\ p95 $<$ 50\,ms};
    \node[note, below=0.55cm of s2] {LLM call, \\ candidate-bounded};
  \end{tikzpicture}
  \caption{Two-stage recommendation pipeline. Stage~1 reads the
    catalog and returns a small candidate set. Stage~2 invokes
    the LLM only over that set and emits identifiers plus a
    rationale. Live product fields are hydrated on the client.
    The dashed arrow indicates that Stage~1 reads from the Convex
  catalog; the LLM never reads the catalog directly.}
  \label{fig:two-stage-pipeline}
\end{figure}

\section{Behavior Tracking and Readiness Signals}
\label{sec:sys-behavior}

Behavior tracking is five client-side hooks in
\texttt{apps/web/src/hooks/}. \texttt{useDwellTime} runs a tick
counter while the cursor is over a product container.
\texttt{useViewportTracking} accumulates
IntersectionObserver-reported visible time.
\texttt{useScrollDepth} maintains the page-level maximum scroll
fraction. \texttt{useProductEngagement} composes the first three
behind a single product-scoped facade.
\texttt{useBehaviorTracker} aggregates everything into the
per-session buffer that flushes to the backend.

Listing~\ref{lst:use-dwell-time} shows the core of
\texttt{useDwellTime}. The hook has two effects: one attaches the
\texttt{mouseenter} and \texttt{mouseleave} listeners, and one
runs a 100~ms tick interval while the cursor is hovering. The
cumulative dwell time is the sum of previous hover sessions plus
the current ongoing tick count, which lets the consuming
component read the live value without waiting for a hover to end.

\begin{listing}[ht]
\begin{minted}{typescript}
useEffect(() => {
  if (!element) return;
  const enter = () => setIsHovered(true);
  const leave = () => {
    setIsHovered(false);
    setCurrentMs((prev) => {
      setCumulativeMs((cum) => cum + prev);
      return 0;
    });
  };
  element.addEventListener("mouseenter", enter);
  element.addEventListener("mouseleave", leave);
  return () => {
    element.removeEventListener("mouseenter", enter);
    element.removeEventListener("mouseleave", leave);
  };
}, [element]);
\end{minted}
  \caption{The hover-tracking effect inside \texttt{useDwellTime}.
    Mouse enter and leave listeners are attached imperatively; on
    leave, the in-progress tick total is folded into the
  cumulative dwell value.}
  \label{lst:use-dwell-time}
\end{listing}

\paragraph{Aggregation and flush.}
\texttt{useBehaviorTracker} keeps a
\texttt{Map<productId, TrackedProduct>} buffer plus a
\texttt{ref} mirror so the flush callback reads current state
without triggering a re-render. Aggregation is lossless within
the flush window (maximum dwell, scroll depth, hover time, and
the boolean OR of \texttt{viewedReviews} per product). Three
subscriptions invoke the flush: a 5-second interval, a
pathname-change effect, and a \texttt{visibilitychange} listener
that catches tab close. Convex performs max-value merging on
conflict, so the same session can flush multiple times without
losing aggregate state (Figure~\ref{fig:behavior-flush}).

\begin{figure}[ht]
  \centering
  \includegraphics[width=0.88\linewidth]{behavior-flush.pdf}
  \caption{Behavior aggregation and flush pipeline. Events from
    four tracking hooks are buffered per product and flushed on
    a 5-second interval, route change, or tab visibility change,
  sending only aggregate fields to the backend.}
  \label{fig:behavior-flush}
\end{figure}

\paragraph{Readiness signals.} Five signals run continuously on
the client. Each maps to an ambient indicator, not an action.
Product dwell longer than 15~s on a detail page (or 5~s on a
card) and deep scroll past 50\% trigger a breathing animation on
the advisor button. Review-section visibility surfaces a
``What do buyers think?'' chip. Three products viewed in a
category without an add-to-cart surfaces a
``Compare these products?'' chip. Items in the cart during
away-navigation surface a ``Need help deciding?'' chip. The
\texttt{useReadinessSignals} hook reads the in-memory behavior
state and returns \texttt{shouldPulseAdvisor} and
\texttt{activeChips}; dismissed chips persist for the session.
A composite interest score over dwell, scroll, hover, and
review-tab engagement is also computed and kept as a
context-richness hint for the advisor prompt, but the simpler
per-signal thresholds proved more legible during early
dogfooding and are what currently drive the indicators. The
evaluator runs entirely client-side; no LLM call is issued by
this loop.

\section{LLM Advisor}
\label{sec:sys-llm}

The advisor is a Convex action wrapping an instance of
\texttt{@convex-dev/agent}, the official Convex agent component.
The component handles thread persistence, message ordering,
streaming, and AI SDK integration. The application code defines
the model selection, the system prompt, and the tools the agent
can call.

\paragraph{Model routing.} Two tiers from a single provider
(OpenRouter) are used. The conversational tier is
\texttt{openai/gpt-oss-120b}, served by Cerebras with Groq as
fallback, at medium reasoning effort. The batch tier is
\texttt{google/gemini-3.1-flash-lite-preview} with reasoning
disabled. Listing~\ref{lst:select-model} shows
\texttt{selectModel} from \texttt{convex/ai/models.ts}: the model
map resolves a query type to a model identifier, and for gpt-oss
models, provider routing prefers Cerebras then Groq for
throughput.

\begin{listing}[ht]
\begin{minted}{typescript}
const FLASH = "openai/gpt-oss-120b";
const FLASH_LITE = "google/gemini-3.1-flash-lite-preview";

const MODEL_MAP: Record<QueryType, string> = {
  conversational: FLASH,
  comparison: FLASH,
  personalized_advice: FLASH,
  review_summary: FLASH_LITE,
  quick_qa: FLASH_LITE,
  homepage_hero: FLASH_LITE,
};

export function selectModel(queryType: QueryType) {
  const modelId = MODEL_MAP[queryType];
  if (modelId.startsWith("openai/gpt-oss")) {
    return openrouter(modelId, {
      reasoning: { effort: "medium" },
      provider: { order: ["cerebras", "groq"], allow_fallbacks: true,
                  sort: "throughput" },
    });
  }
  return openrouter(modelId);
}
\end{minted}
  \caption{The model-routing function. The \texttt{provider.order}
    array pins Cerebras as the preferred backend for gpt-oss-120b;
    \texttt{allow\_fallbacks} permits OpenRouter to switch to Groq
    if Cerebras is unavailable. The throughput sort matches the
  latency-sensitive nature of the conversational tier.}
  \label{lst:select-model}
\end{listing}

\paragraph{Agent, tools, and prompt.} The agent is constructed
once per cold start as a module-level constant with
\texttt{maxSteps: 12} as the step budget. The value is tuned by
hand: \texttt{gpt-oss-120b} is aggressive with tool calls (three
to five per turn is typical), and an earlier value of 5 caused
the agent to exhaust the budget before producing a final text
answer. That is the failure the eval harness's
\texttt{hasFinalAnswer} scorer was added to catch automatically.
Five tools are registered: \texttt{getProductDetails},
\texttt{searchProducts}, \texttt{getRecommendations} (which
dispatches to similar, frequently-bought-together, or trending),
\texttt{getCartContents}, and \texttt{getReviewsSummary} (which
returns the structured summary including the conflicts list).
Each tool wraps a Zod input schema and an async execute function
via the \texttt{createTool} helper. No tool generates free text
and no tool returns synthesized data: every field is read
directly from the live Convex catalog. This is
hydration-over-generation in action
(Section~\ref{sec:found-grounding}). The system prompt encodes
the design constraints as hard rules (``always use tool results
for factual claims, never guess prices, ratings, or
specifications''; ``never recommend products already in the
customer's cart''; ``when discussing reviews, surface both
positives and negatives with counts; if conflicts exist, always
mention them''). The prompt is the first line of defense; the
validation pipeline in Section~\ref{sec:sys-validation} is the
structural guarantee.

\paragraph{Advisor sidebar UI.} The advisor surface is a
push-content sidebar (page content shifts left instead of being
overlaid by a modal). Four components compose it under
\texttt{components/advisor/}: the floating button with the
breathing animation, the sidebar container, the input composer,
and the message list. Streaming comes from the agent component's
\texttt{useUIMessages} hook, keyed on a thread ID created with a
separate \texttt{createThread} mutation so the client subscribes
before the streaming action begins.
Figure~\ref{fig:advisor-request} traces a single end-to-end
advisor request.

\begin{figure}[ht]
  \centering
  \includegraphics[width=0.92\linewidth]{advisor-request.pdf}
  \caption{End-to-end advisor request. The user click is the
    only entry point. The client assembles context from the
    readiness and behavior hooks. The Convex action invokes the
    agent with a model routed per query type. The agent runs a
    tool-calling loop against the catalog. Validation strips any
    unverifiable claims. The client hydrates the productIds
    against the live catalog and renders. No factual data flows
  from the LLM to the UI directly.}
  \label{fig:advisor-request}
\end{figure}

\section{Review Summarization}
\label{sec:sys-reviews}

Per-product review summaries serve two surfaces: the product
detail page and the advisor. The summary is generated offline,
batch-style, against the actual review corpus. No summary is
ever generated on a user-facing critical path. The structure is
fixed: \emph{what buyers like}, \emph{common complaints},
\emph{best for}, and \emph{conflicting opinions} (explicit
disagreement pairs, each annotated with per-side counts). The
conflict-surfacing structure is the anti-Rufus design choice
from Chapter~\ref{chap:background}: ``45 reviews praise battery
life, 12 reviews call it disappointing'' is more useful than
``most users are happy with the battery''.

The generation flow is a daily Convex cron that finds products
whose summary is stale (missing, or whose review count differs
from the current count, or older than seven days) and
regenerates each. Generation reads the full review corpus,
batches it into a Flash-Lite prompt with a structured JSON
schema, and persists the result; every theme carries a
representative quote and the originating reviewId. After
generation, each theme is verified against the review corpus
using the embedding index on the reviews table: the verifier
embeds the theme text, runs a vector search scoped to the
product, and accepts the theme only if at least one review in
the top-$k$ neighbors clears the cosine-similarity threshold
$\tau$ from Section~\ref{sec:found-grounding}. Failing themes
are dropped before persistence; the \texttt{validationDetails}
field records how many were verified versus dropped.

\section{AI-Assisted Comparison}
\label{sec:sys-comparison}

The comparison flow has no Convex action or agent of its own. It
reuses the advisor entry point with two changes. First, the
request payload includes a \texttt{queryType: "comparison"}
hint, which the action uses to pick the comparison-specific
system prompt variant. Second, the agent is invoked with a
smaller step budget (the comparison flow does not need the
tool-calling depth a freeform conversation does), with the same
five tools registered.

The comparison-specific prompt variant restricts the answer to a
structured tradeoff format, forbids a single-winner verdict, and
refuses comparisons across incompatible categories or duplicate
SKUs. Three structural choices define the surface: the user
picks the products (the system does not), the comparison is
multidimensional across price, specs, ratings, and review
themes, and no winner is named. The user makes the call. This is
a deliberate departure from how most shopping AI handles
comparison; the related-work survey in
Section~\ref{sec:bg-comparison} arrives at the same
recommendation~\cite{ShopifyAgentic,SmashingAgenticUX}. The
structured output renders as a side-by-side table built from the
live catalog using the LLM-emitted productIds, plus a free-text
tradeoff narrative; validation (Section~\ref{sec:sys-validation})
strips price and specification claims from the narrative.

\section{Output Validation Pipeline}
\label{sec:sys-validation}

Output validation is implemented at the boundary between the
agent action and the client response, inside the same Convex
action that wraps the agent invocation. The three validators run
sequentially as formalized in
Section~\ref{sec:found-grounding}. Figure~\ref{fig:validation-pipeline}
shows the pipeline.

\begin{figure}[ht]
  \centering
  \includegraphics[width=0.94\linewidth]{validation-pipeline.pdf}
  \caption{Output validation pipeline. The structured part of the
    LLM response is checked against the live catalog (productId
    existence). The free-text rationale is scrubbed of price and
    specification claims. Review-summary themes are verified
    against the review corpus before persistence. Failures
  result in dropped references rather than user-visible errors.}
  \label{fig:validation-pipeline}
\end{figure}

\paragraph{Identifier-existence.} Every productId reference the
agent emits is checked against the catalog via a single
\texttt{ctx.runQuery(api.products.get, ...)} per identifier.
Identifiers in $\hat{S} \setminus \mathcal{C}_{id}$ are stripped
from the response before the action returns. The cost is one
indexed read per cited identifier; in practice an advisor turn
cites at most a handful of products, so the latency overhead is
sub-millisecond.

\paragraph{Claim stripping.} The free-text rationale runs
through a small set of regular expressions that match
currency-prefixed numbers, decimal patterns next to keywords
like ``costs'' or ``priced at'', and unit-suffixed numbers near
keywords like ``GB'', ``inches'', ``mAh'', or ``Hz''. Matches
are removed from the rationale. The system prompt already
forbids these claims; the regex pass is the second-line defense
for prompt-following lapses.

\paragraph{Theme verification.} Verification only applies to
outputs that include claimed review themes (the review-summary
surface and certain advisor responses that quote review
feedback). The verifier reuses the embedding-search machinery
from Section~\ref{sec:sys-reviews}: each claimed theme is
embedded, a vector search runs against reviews scoped to the
relevant product, and the theme is dropped if no review in the
top-$k$ clears the similarity threshold.

The three checks are not redundant. Identifier-existence catches
structural hallucinations. Stripping catches drift in free text.
Theme verification catches fabrication in structured outputs
that pass the schema check but are not supported by source data.
Together they encode the design constraint that the LLM is a
candidate selector and a writer, never a source of facts.

\section{Frontend Walkthrough}
\label{sec:sys-frontend}

Two surfaces illustrate the visual contract the system makes
with the user. The first is the homepage with its recommendation
rows. The second is the AI-generated review summary on the
product detail page.

\begin{figure}[ht]
  \centering
  \includegraphics[width=0.95\linewidth]{ui-homepage-recommendations.png}
  \caption{Homepage surface. Two consecutive recommendation rows
    are visible: a popularity-driven row at the top and the
    personalized ``Recommended for you'' row below it. The
    floating advisor button in the lower-right corner is always
  present and idle until a readiness signal fires.}
  \label{fig:ui-homepage}
\end{figure}

The homepage is composed entirely of classical recommendation
slots: a popularity row on top, the personalized
``Recommended for you'' switching-hybrid output below. None of
these invokes an LLM call. The entire homepage renders from
Stage~1. The advisor button in the lower-right corner is idle
until a readiness signal fires, expressing the reactive-first
model visually: the AI is available, never imposed.

\begin{figure}[ht]
  \centering
  \includegraphics[width=0.95\linewidth]{ui-product-review-summary.png}
  \caption{AI-generated review summary on the product detail
    page. The bipartite conflicts bar surfaces divided opinions
    visually, with per-side counts and a one-sentence
    characterization of each side. The ``Best for'' pills
    surface the structured buyer-cohort field from the summary
  schema. A small ``AI-generated'' tag discloses the synthetic origin.}
  \label{fig:ui-review-summary}
\end{figure}

The review-summary surface concentrates the anti-Rufus design
choices. The explicit positives-and-complaints split shows each
theme with a per-theme count and a representative quote. The
bipartite conflicts bar visualizes divided opinions with
agreeing- and disagreeing-review counts. The ``Best for'' pills
expose the buyer-cohort field. A small ``AI-generated'' tag
discloses the synthetic origin. The contextual
``What do buyers think?'' chip is the review-engagement
readiness signal: it dispatches a question to the advisor on
click but never triggers an LLM call by itself. Neither surface
displays a number, attribute, or price that originates in the
LLM rather than in the live catalog.

\section{Evaluation Harness}
\label{sec:sys-eval}

The eval harness lives in \texttt{packages/eval} and is built
around Promptfoo as the runner with custom shopping-domain
scorers. Section~\ref{sec:found-eval} formalized the scorers and
the composite ranking; this section describes the implementation.
Promptfoo runs the cartesian product of dataset rows and provider
configurations. A custom \texttt{advisorProvider} invokes a
Convex action wrapping the production advisor agent, and
programmatic, built-in, and LLM-as-judge scorers all write into a
single JSON results file per sweep.

\paragraph{Build vs.\ buy.} An honest accounting of where the
engineering effort has thesis value left only the scoring as
distinctive: whether an answer cited a real product, whether a
quoted price matches the database, whether a review theme is
supported by the actual review corpus. These scorers depend on
the live catalog and the embedded review index, and they do not
exist in any off-the-shelf tool. The implementation hands the
generic parts (dataset loading, sweep orchestration, result
storage, local dashboard) to Promptfoo (MIT, open-source) and
ships the shopping-domain scorers as custom JavaScript
assertions.

\paragraph{Dataset.} The dataset is
\texttt{src/datasets/shopping-v1.yaml}, a 25-row YAML file with
each row carrying the question, a simulated context
(\texttt{productId}, \texttt{recentlyViewedIds}), the expected
tools the agent should call, and the productIds the answer
should reference. Rows are distributed across seven categories
(Table~\ref{tab:eval-dataset-categories}). Dataset changes are
versioned by filename, so historical runs are never compared
against new questions.

\begin{table}[ht]
  \centering
  \footnotesize
  \begin{tabular}{|l|c|l|}
    \hline
    \textbf{Category} & \textbf{Rows} & \textbf{What it stresses} \\
    \hline
    simple Q\&A             & 4 & single-tool factual lookup \\
    product validation      & 4 & confirming a single product detail \\
    comparison              & 5 & multi-product tradeoff structure \\
    recommendation          & 5 & switching hybrid + rationale \\
    review summary          & 3 & grounded use of summarized reviews \\
    multi-turn              & 3 & context carryover across turns \\
    edge case               & 1 & ambiguous, off-topic, empty-result \\
    \hline
    \textbf{Total}          & \textbf{25} & \\
    \hline
  \end{tabular}
  \caption{Distribution of dataset rows in
  \texttt{shopping-v1.yaml} across seven categories.}
  \label{tab:eval-dataset-categories}
\end{table}

\paragraph{Configuration sweep.} Nine configurations are
evaluated in the canonical sweep
(Table~\ref{tab:eval-sweep-configs}). The sweep varies four axes:
model identifier, reasoning effort, step budget, and prompt
variant. Promptfoo runs the cartesian product
($9 \times 25 = 225$ advisor invocations).

\begin{table}[ht]
  \centering
  \footnotesize
  \begin{tabular}{|l|l|c|c|l|}
    \hline
    \textbf{Label} & \textbf{Model} & \textbf{Reasoning} &
    \textbf{Steps} & \textbf{Prompt variant} \\
    \hline
    baseline-flash-lite      & gemini-3.1-flash-lite-preview & n/a    & 12 & current \\
    gpt-oss-120b-low         & gpt-oss-120b                  & low    & 12 & current \\
    gpt-oss-120b-medium      & gpt-oss-120b                  & medium & 12 & current \\
    gpt-oss-120b-high        & gpt-oss-120b                  & high   & 12 & current \\
    gpt-oss-20b-low          & gpt-oss-20b                   & low    & 12 & current \\
    gpt-oss-20b-medium       & gpt-oss-20b                   & medium & 12 & current \\
    oss-120b-tight           & gpt-oss-120b                  & medium & 8  & ``be efficient'' \\
    oss-120b-loose           & gpt-oss-120b                  & medium & 15 & current \\
    oss-120b-no-fewshot      & gpt-oss-120b                  & medium & 12 & no few-shot \\
    \hline
  \end{tabular}
  \caption{Configuration sweep for the canonical evaluation.
    Each configuration is run against all 25 dataset rows,
    producing 225 advisor invocations and approximately 1\,125
  LLM-judge calls.}
  \label{tab:eval-sweep-configs}
\end{table}

\paragraph{Scorers.} Five programmatic scorers
(\texttt{hasFinalAnswer}, \texttt{groundedness},
  \texttt{factuality}, \texttt{reviewThemeFidelity},
\texttt{toolCallEfficiency}) plus Promptfoo's built-in
\texttt{tool-call-f1} are registered. Listing~\ref{lst:groundedness}
shows \texttt{groundedness.ts} in full. It is the scorer most
directly aimed at the Rufus 32\% accuracy failure mode.

\begin{listing}[ht]
\begin{minted}{typescript}
const PRODUCT_ID_REGEX = /\bk[0-9a-z]{31}\b/g;

export default async function groundedness(output: string, _ctx: Context) {
  const cited = Array.from(new Set(output.match(PRODUCT_ID_REGEX) ?? []));
  if (cited.length === 0) {
    return { pass: true, score: 1, reason: "vacuously grounded" };
  }
  const lookups = await Promise.all(cited.map(async (id) => {
    try {
      const product = await convex.query(productGetRef, { id });
      return { id, exists: product !== null };
    } catch { return { id, exists: false }; }
  }));
  const missing = lookups.filter((l) => !l.exists).map((l) => l.id);
  const grounded = cited.length - missing.length;
  return {
    pass: missing.length === 0,
    score: grounded / cited.length,
    reason: missing.length === 0
      ? `all ${cited.length} cited productIds exist`
      : `${missing.length}/${cited.length} hallucinated: ${missing.join(", ")}`,
  };
}
\end{minted}
  \caption{The \texttt{groundedness} scorer. Convex IDs are
    32-character URL-safe alphanumerics; products start with a
    lowercase ``k''. Every cited identifier is looked up against
    the live catalog through a typed Convex client. A perfect
  score requires every cited identifier to exist.}
  \label{lst:groundedness}
\end{listing}

Five LLM-as-judge rubrics (completeness, helpfulness, tradeoff
surfacing, tool appropriateness, tone) run on a 0--5 scale with
Claude Haiku~4.5 as the judge model, chosen across model families
to mitigate self-preference bias as formalized in
Section~\ref{sec:found-eval}.

\section{Results}
\label{sec:sys-results}

This section reports the measured findings from the canonical
sweep. The composite score combines judge quality (weight 0.30),
programmatic correctness (0.25), per-query dollar cost (0.20),
p95 latency (0.15), and tool-step efficiency (0.10); cost and
latency enter with negative weight and a per-dimension normalizer
so the composite stays in $[0,1]$. Every configuration ran the
full 34 rows of the dataset (25 questions plus four-rubric judge
invocations), so row coverage is uniform across cells.

\subsection{Ranked Configurations}
\label{subsec:results-sweep}

Table~\ref{tab:results-ranked-configs} reports the ranking.

\begin{table}[ht]
  \centering
  \footnotesize
  % generated by packages/eval/src/reports/thesisExport.ts — do not edit by hand
\begin{tabular}{rlrrrrrr}
\toprule
\# & configuration & quality & correctness & efficiency & avg cost & p95 latency (ms) & composite \\
\midrule
1 & gemini-3.1-flash-lite-preview · none · ms12 · current & 0.776 & 0.941 & 0.515 & \$0.0003 & 4216 & \textbf{0.807} \\
2 & gpt-oss-120b · low · ms12 · current & 0.775 & 0.922 & 0.515 & \$0.0006 & 4705 & \textbf{0.795} \\
3 & gpt-oss-120b · medium · ms12 · no-fewshot & 0.815 & 0.921 & 0.497 & \$0.0006 & 5889 & \textbf{0.774} \\
4 & gpt-oss-20b · medium · ms12 · current & 0.672 & 0.920 & 0.506 & \$0.0005 & 4181 & \textbf{0.768} \\
5 & gpt-oss-120b · medium · ms8 · be-efficient & 0.703 & 0.932 & 0.497 & \$0.0007 & 5734 & \textbf{0.719} \\
6 & gpt-oss-120b · medium · ms15 · current & 0.726 & 0.930 & 0.479 & \$0.0007 & 7466 & \textbf{0.683} \\
7 & gpt-oss-120b · medium · ms12 · current & 0.764 & 0.924 & 0.485 & \$0.0007 & 6053 & \textbf{0.680} \\
8 & gpt-oss-120b · high · ms12 · current & 0.734 & 0.900 & 0.467 & \$0.0010 & 12072 & \textbf{0.492} \\
\bottomrule
\end{tabular}

  \caption{Configurations ranked by composite score. Each row is
    a single configuration evaluated against the 25 rows of
  \texttt{shopping-v1.yaml}.}
  \label{tab:results-ranked-configs}
\end{table}

Three observations follow. The Gemini~3.1 Flash-Lite baseline
ranks \#1 on composite (0.807) despite being treated by the
original design as a cheaper fallback. It loses on raw quality
by 3.9 points to \texttt{no-fewshot} \texttt{gpt-oss-120b}
(0.776 vs.\ 0.815) but wins decisively on both penalty axes:
\$0.0003 per query versus \$0.0005--0.0010, and 4.2~s p95
versus 4.2--12.1~s. Second, \texttt{gpt-oss-120b} at \emph{low}
reasoning ranks \#2 (0.795), and \texttt{no-fewshot} ranks \#3
(0.774) while still producing the highest raw quality in the
sweep. Two things follow: the few-shot examples in the
production prompt are a net negative on this dataset
(Section~\ref{subsec:results-prompt-variant}), and the default
\texttt{medium} reasoning setting is not the composite-optimal
one. Third, raising reasoning effort and loosening the step
budget both hurt the composite uniformly:
\texttt{gpt-oss-120b} high-reasoning is the worst at 0.492
despite competitive raw quality (cost doubles, p95 inflates to
12.1~s), and \texttt{ms15} comes in at 0.683 with raw quality
\emph{below} the \texttt{ms12} default (the extra steps did not
improve the answer).

\subsection{Pareto Fronts}
\label{subsec:results-pareto}

Figures~\ref{fig:results-pareto-cost-quality}
and~\ref{fig:results-pareto-latency-quality} plot each sweep cell
on the (cost, quality) and (latency, quality) planes.

\begin{figure}[ht]
  \centering
  \input{figures/eval-pareto-cost-quality}
  \caption{Cost-quality pareto plot. $x$-axis is average
    per-query cost in USD, $y$-axis is the LLM-judge mean
    quality. Flash-Lite (leftmost) and \texttt{no-fewshot}
  (topmost) are the two non-dominated configurations.}
  \label{fig:results-pareto-cost-quality}
\end{figure}

\begin{figure}[ht]
  \centering
  % generated by thesisExport.ts — do not edit by hand
% requires: \usepackage{pgfplots}
\begin{tikzpicture}
  \begin{axis}[
    width=\linewidth,
    height=8cm,
    xlabel={Average latency (ms)},
    ylabel={Quality (LLM-judge mean)},
    grid=both,
    nodes near coords,
    nodes near coords align={vertical},
    every node near coord/.append style={font=\tiny}
  ]
    \addplot+[only marks, mark=*, point meta=explicit symbolic]
      table[col sep=comma, x=latency, y=quality, meta=label] {eval-pareto-latency-quality.csv};
  \end{axis}
\end{tikzpicture}

  \caption{Latency-quality pareto plot. Flash-Lite anchors the
    low-latency end at competitive quality; \texttt{no-fewshot}
  anchors the high-quality end at moderately higher latency.}
  \label{fig:results-pareto-latency-quality}
\end{figure}

On the cost-quality plane, only two configurations sit on the
frontier. Flash-Lite is unbeaten on cost at the price of 3.9
quality points; \texttt{no-fewshot} is unbeaten on quality at
the price of roughly doubling per-query cost. Every other
configuration lies in the strictly-dominated interior. The
implication for production is that the choice between
Flash-Lite and \texttt{no-fewshot} is the only choice that has
to be made on this evidence. The latency-quality plane tells
the same structural story. Flash-Lite and \texttt{no-fewshot}
again anchor the frontier, with Flash-Lite cheaper \emph{and}
faster (4.2~s p95 versus 5.9~s) at the cost of those 3.9
quality points.

\subsection{Per-Category Programmatic Correctness}
\label{subsec:results-per-category}

Table~\ref{tab:results-per-category} reports per-cell
programmatic-correctness means (averaged across the six
deterministic scorers) restricted to the three categories with
enough rows to support a breakdown.

\begin{table}[ht]
  \centering
  \footnotesize
  % generated by thesisExport.ts — do not edit by hand
\begin{tabular}{lrrrrrr}
\toprule
configuration & comparison & edge\_case & product\_validation & recommendation & review\_summary & simple\_qa \\
\midrule
gemini-3.1-flash-lite-preview · none · ms12 · current & 0.95 & 1.00 & 0.97 & 0.89 & 0.93 & 0.98 \\
gpt-oss-120b · low · ms12 · current & 0.91 & 1.00 & 0.95 & 0.84 & 1.00 & 1.00 \\
gpt-oss-120b · medium · ms12 · no-fewshot & 0.91 & 1.00 & 0.94 & 0.81 & 1.00 & 1.00 \\
gpt-oss-20b · medium · ms12 · current & 0.94 & 1.00 & 0.91 & 0.85 & 1.00 & 0.97 \\
gpt-oss-120b · medium · ms8 · be-efficient & 0.92 & 1.00 & 0.95 & 0.90 & 0.97 & 0.96 \\
gpt-oss-120b · medium · ms15 · current & 0.91 & 1.00 & 0.94 & 0.82 & 0.97 & 1.00 \\
gpt-oss-120b · medium · ms12 · current & 0.93 & 1.00 & 0.91 & 0.83 & 0.97 & 0.99 \\
gpt-oss-120b · high · ms12 · current & 0.92 & 1.00 & 0.89 & 0.79 & 0.85 & 0.94 \\
\bottomrule
\end{tabular}

  \caption{Mean programmatic correctness per category, per
    configuration. Values are bounded in $[0,1]$, higher is
  better.}
  \label{tab:results-per-category}
\end{table}

Programmatic correctness is uniformly high on \texttt{simple\_qa}
(0.94--1.00): when the question is a single factual lookup,
every configuration calls the right tool, the right product
exists, and the price and rating match. This is the floor the
catalog-hydration design is meant to enforce, and the harness
confirms it. \texttt{Recommendation} is where configurations
actually differ, ranging from 0.79 (high-reasoning) to 0.90
(\texttt{ms8} \texttt{be-efficient}). That category is the
structurally hardest case (open-ended candidate ranking,
tradeoff language, no single ground-truth product), and the one
where future dataset growth would produce the most discriminating
signal.

\subsection{Few-Shot Contamination}
\label{subsec:results-prompt-variant}

The harness was built to surface failure modes that the
production UI hides. The cleanest example is the few-shot
contamination the prompt-variant A/B revealed. With the
production prompt (\texttt{gpt-oss-120b}, medium reasoning,
12-step budget, few-shot examples included), the advisor's
responses on simple, ambiguous questions systematically
reference products from the few-shot examples (Logitech~K380
keyboard, Samsung Galaxy~S24~Ultra, iPhone~15~Pro) instead of
the product carried in the row's system context. Removing the
few-shot examples eliminates the contamination.

The qualitative form is clearest on a single row. The dataset
row asks for the price of a Vivo~V9; the row's system context
names the product explicitly.

\begin{quote}
\textbf{\texttt{current} (few-shots included).} ``I'm happy to
look up the price and any current promotions, but could you let
me know which phone you're referring to? \emph{Samsung Galaxy
S24 Ultra}, \emph{iPhone 15 Pro}, or one of the other models
you've been browsing\ldots''

\medskip

\textbf{\texttt{no-fewshot}.} ``The Vivo V9 is priced at
\$299.99. At the moment there's no discount listed for this
model. If a promotion becomes available, it would be reflected
on the product page.''
\end{quote}

The mechanism is mechanical: few-shot examples are injected into
the model's \texttt{messages} array as JSON-encoded
\texttt{user} and \texttt{assistant} turns. On a short ambiguous
query like ``what's the price of this phone?'' the model treats
those turns as recent conversation history and pattern-matches
against the example products instead of the system-context
product. The production UI largely side-steps the failure
because real users arrive with specific questions (from
comparison chips, review-summary buttons), but the harness's
broader query distribution surfaces it. Programmatic and judge
measurements agree on direction: \texttt{no-fewshot} hits the
highest raw judge quality (0.815 vs.\ 0.764 for the
otherwise-identical \texttt{current}) while keeping correctness
within 0.3~pp, and a focused two-configuration A/B from earlier
development ranked the variants identically on all four judge
rubrics. Two scoring methodologies converging on the same
ordinal is what argues the finding is not an artifact of either.

\subsection{Summary and Validity}
\label{subsec:results-summary}

The evidence supports the central claim on two of three axes:
every configuration completes well under the
consumer-tolerance latency bound and at well under one cent per
query, with grounded output across the sweep
(\texttt{groundedness} mean $\ge 0.89$). The harness adds a
finding the claim did not anticipate: a non-trivial design
space exists \emph{within} those bounds where the right choice
is cheaper and less computationally aggressive rather than the
one with the highest raw quality. Higher reasoning effort and
looser step budgets are uniformly worse on composite, which
disproves the implicit ``more reasoning is always better''
assumption. Three validity classes qualify these numbers.
\emph{Internal}: seeded reviews are AI-generated, and residual
judge bias remains despite the cross-family judge choice.
\emph{Construct}: groundedness verifies existence not relevance,
factuality verifies values not semantic correctness, theme
fidelity verifies embedding match not salience. The composite
is reported alongside the pareto fronts, so disagreement with
the weights is recoverable. \emph{External}: 200 products on a
development-tier deployment do not predict 10\,000-product
production load, and the 25-row dataset cannot cover every
failure mode. The methodological finding (programmatic and
cross-family judges agree on ordinal) argues the harness
produces a robust ranking even when individual scorers are
imperfect.

\section{Engineering Challenges}
\label{sec:sys-engineering}

Six engineering challenges hit during implementation justify
choices that affect the metrics above. Each is documented at
length in \texttt{docs/lessons-learned.md}.

\begin{itemize}
  \item \textbf{Hooks before returns.} Every React hook appears
    before any conditional control flow; rules-of-hooks
    violations are caught by strict mode and the React Compiler.
  \item \textbf{Hydration mismatches with Base UI.} Components
    are gated by a \texttt{useMounted} hook implemented with
    \texttt{useSyncExternalStore} (\texttt{getServerSnapshot:
      () => false}), collapsing the typical two-render pattern
    into one.
  \item \textbf{Convex action abort.}
    \texttt{useAction} does not support abort signals; the
    advisor stop button clears the UI loading state but the
    server-side action runs to completion (cost acknowledged in
    the per-turn token figures).
  \item \textbf{\texttt{useUIMessages} typing.}
    \texttt{UIMessagesQueryResult<typeof
      api.ai.queries.listMessages>} derives the message type
    from the source of truth instead of duplicating it.
  \item \textbf{\texttt{content-visibility} grid
    virtualization.} CSS-native
    \texttt{content-visibility: auto} with
    \texttt{contain-intrinsic-size} for 50--150-card grids,
    avoiding any JavaScript virtualizer.
  \item \textbf{100~ms tracking-hook intervals.} Per-element
    intervals are cheap on desktop hardware; potentially
    measurable on low-end mobile.
\end{itemize}
